% Evans 11.4, printed pp.601--613; English PDF pp.618--630.
% Chinese start: printed p.425 / scan p.443.
\section{Entropy Criteria}
\label{sec:11-4}
\subsection{Viscous profiles and the limits of dimension counting}
For identity viscosity,
$u_t^\varepsilon+F(u^\varepsilon)_x=\varepsilon u_{xx}^\varepsilon$.
The ansatz $u^\varepsilon=V((x-\sigma t)/\varepsilon)$ gives
\begin{equation}\label{eq:11-profile}
 \begin{split}
 V'&=F(V)-F(u_l)-\sigma(V-u_l)=G(V;\sigma),\\
 V(-\infty)&=u_l,\qquad V(+\infty)=u_r.
 \end{split}
\end{equation}
Here integration used $V'(-\infty)=0$; the right limit then forces
Rankine--Hugoniot. For an autonomous smooth ODE, convergence of $V$
to an equilibrium also implies $V'\to0$.
\PDEAnnotation{稳定维数不等于轨道存在}{线性化特征值是
$\lambda_j(u_l)-\sigma$。某个方向不稳定只说明允许离开平衡点，
不能单凭维数判断一条轨道必到达指定的右状态；一般系统须另用中心流形等工具。}

We use the following \emph{external finite-dimensional theorem}, not
proved here: a sufficiently smooth autonomous vector field with a
spectral splitting into center, stable, and unstable subspaces has a
local center manifold of any fixed finite differentiability allowed
by the field. It is tangent to the center space and contains every
complete orbit staying in a sufficiently small neighborhood, in
particular all nearby equilibria. Parameters can be included as
variables with zero derivative. The conclusion is local; it does not
assert that a distant excursion lies on the manifold.

\begin{theorem}[Small local viscous profiles]
This is a local version of source \S11.4.1, Theorem 1.
Assume $F$ is smooth and the $k$-field is genuinely nonlinear near
$u_l$. For a sufficiently small nonzero $k$-shock state $u_r$, a
profile whose entire image stays near $u_l$ exists if and only if
it is on the compressive branch $S_k^-(u_l)$. The proof below uses
the stated external center-manifold theorem.
\end{theorem}
\begin{expandedproof}[center-manifold reduction]
Fix $u_l$, orient $r=r_k(u_l)$ so $\gamma=D\lambda_k(u_l)r>0$,
and write $\beta=\sigma-\lambda_k(u_l)$. Adjoin $\beta'=0$ to
\eqref{eq:11-profile}. At $(V,\beta)=(u_l,0)$ the center space has
dimension two, spanned by $(r,0)$ and $(0,1)$; all other eigenvalues
are separated from zero by strict hyperbolicity. A local center
manifold can be parameterized by $a=l_k(u_l)(V-u_l)$ and $\beta$:
\[
 V=u_l+ar+W(a,\beta),\quad l_kW=0,\quad
 W(0,\beta)=0,\quad W=O(a^2+|a\beta|).
\]
The equality $W(0,\beta)=0$ follows because $V=u_l$ is an equilibrium
for every small $\beta$. Project the vector field by $l_k$.
Taylor expansion, $l_k(DF-\lambda_kI)=0$, and $l_kr=1$ give
\[
 a'=g(a,\beta)=-\beta a+\tfrac12\gamma a^2
       +O\bigl(|a|(|a|+|\beta|)^2\bigr).
\]
The coefficient identity $l_kD^2F[r,r]=\gamma$ follows by
differentiating the eigenvalue equation. The nearby right equilibrium
lies on the manifold. Write its coordinate $a_r\ne0$; by the shock
speed expansion, $\beta=\gamma a_r/2+O(a_r^2)$.
For this fixed $\beta$, both $0$ and $a_r$ are roots of $g$.
Divide smoothly first by $a$ and then by $a-a_r$, using the integral
form of the mean value theorem. Uniformly for $a$ between the roots,
\[
 g(a,\beta)=a(a-a_r)h(a,\beta),\qquad
 h=\gamma/2+O(|a_r|)>0.
\]
If $a_r<0$, the vector field is negative between $a_r$ and $0$.
The solution through any interior point decreases, stays between
the equilibria, and converges to $0$ as $s\to-\infty$ and to $a_r$
as $s\to+\infty$. The simple roots make the times to both endpoints
infinite (the separated integral has logarithmic divergences).
Lifting gives the desired profile, with its full image small.
If $a_r>0$, the field is negative on $(0,a_r)$ and cannot connect
$0$ to $a_r$ forward in $s$. Every complete sufficiently small
profile lies on the center manifold, so this also proves the local
necessity. The equivalence $a_r<0\iff u_r\in S_k^-(u_l)$ follows
from \eqref{eq:11-shock-expansion}. No assertion about profiles
making large excursions is obtained from this local argument.
\end{expandedproof}

For a non-genuinely nonlinear branch, the Liu criterion compares
shock speeds along the full intervening Hugoniot arc:
\[
 \sigma(z,u_l)\ge\sigma(u_r,u_l)
 \quad\hbox{for intermediate states }z\in S_k(u_l).
\]
The source writes a strict version for nontrivial intermediate states.
The local general existence theory requires additional hypotheses and
is not proved by the preceding genuinely nonlinear reduction.
In the scalar profile calculation below, strictness has a precise
meaning: it excludes an extra equilibrium that would block passage.

\subsection{A complete profile test for the viscous elastic system}
Use viscosity only in the velocity equation:
$a_t-b_x=0$, $b_t-p(a)_x=\varepsilon b_{xx}$.
For end strains $a_l<a_r$ choose
\[
 \sigma=\sqrt{\frac{p(a_r)-p(a_l)}{a_r-a_l}}>0,\qquad
 b_r=b_l-\sigma(a_r-a_l),\qquad p'>0.
\]
\begin{theorem}[Source \S11.4.1, elastic example]
A traveling profile joining these states exists if and only if
\begin{equation}\label{eq:11-p-liu}
 p(a)-p(a_l)>\sigma^2(a-a_l),\qquad a_l<a<a_r.
\end{equation}
It is unique up to translation among such monotone profiles.
\end{theorem}
\begin{expandedproof}[separation and infinite endpoint time]
The first integrated equation gives $B=b_l-\sigma(A-a_l)$.
The second yields
\[
 A'=g(A)=\frac{p(A)-p(a_l)-\sigma^2(A-a_l)}\sigma.
\]
Thus $g(a_l)=g(a_r)=0$. If $g>0$ in the open interval, the inverse
of $s-s_0=\int_{A_0}^{A}d\zeta/g(\zeta)$ is a strictly increasing
profile. Since $g$ is $C^1$ and vanishes at each endpoint,
$|g(a)|\le C|a-a_l|$ near $a_l$ and similarly near $a_r$;
the endpoint integrals diverge, even if the roots are not simple.
The inverse is defined for all real $s$ with the required limits.
Conversely any connecting scalar orbit must pass every intermediate
value. Uniqueness for a scalar Lipschitz ODE prevents passage through
an interior zero and prevents a change of direction. Hence $g>0$
throughout the interval. Separation also proves translation uniqueness.
Condition \eqref{eq:11-p-liu} is precisely strict decrease of the
secant speed from $a_l$ to an intermediate point relative to the
final positive speed. This proof concerns the stated partial viscosity,
not an unmentioned change of viscosity matrix.
\end{expandedproof}

\subsection{Entropy pairs, symmetrizers, and viscous limits}
\begin{definition}[Source \S11.4.2]
An entropy pair $(\eta,Q)$ consists of $C^2$ functions with $\eta$
convex and $DQ=D\eta\,DF$. An integral solution is an entropy solution
if for every such pair and $0\le\varphi\in C_c^\infty(\R\times(0,\infty))$,
\begin{equation}\label{eq:11-entropy}
 \iint[\eta(u)\varphi_t+Q(u)\varphi_x]\,dx\,dt\ge0.
\end{equation}
Thus the distribution $\eta(u)_t+Q(u)_x$ is nonpositive.
\end{definition}
For a smooth solution the entropy balance is equality. At a shock it
reduces to $[Q]-\sigma[\eta]\le0$ by the same interface calculation
as before. The sign refers to convex mathematical entropy, often
the negative of physical entropy.
\PDEAnnotation{系统熵对的相容性}{标量情形每个凸 $\eta$ 都能积分得到
$Q$；系统须满足超定条件 $D^2\eta\,DF$ 对称。并不是任意凸函数
都可充当某个系统的熵，更不能由一个能量不等式直接得到所有解的唯一性。}

\begin{lemma}[Entropy compatibility; editorial supplement]
Locally, $D\eta\,DF$ is the gradient of a flux if and only if
$D^2\eta\,DF$ is symmetric. If $D^2\eta$ is positive definite,
$DF$ is similar to a symmetric matrix and has real eigenvalues.
This does not force distinctness.
\end{lemma}
\begin{expandedproof}[integrability and symmetrization]
For $b_j=\sum_a\eta_a F^a_j$, differentiate with respect to $z_i$.
In $\partial_i b_j-\partial_j b_i$ the terms involving $F^a_{ij}$
cancel. The remaining expression is the skew part of $D^2\eta DF$.
A closed $C^1$ one-form has a potential on a small ball: integrate
along radial segments and differentiate, using equality of the cross
partials. This proves local sufficiency; global sufficiency requires,
for example, a simply connected domain. If $S=D^2\eta>0$, then
$S^{1/2}DF S^{-1/2}$ is symmetric because $SDF=(DF)^TS$.
It is similar to $DF$ and hence has the same real spectrum.
\end{expandedproof}
For the elastic system, a strictly convex pair is
$\eta(a,b)=b^2/2+P(a)$, $P'=p$, and $Q=-bp(a)$.
Indeed $D^2\eta=\operatorname{diag}(p',1)>0$.

\begin{theorem}[Viscous limit; source \S11.4.2, Theorem 2]
Suppose smooth identity-viscosity solutions with initial data $g$
are uniformly bounded with range in a fixed compact state set and
converge a.e. to $u$ on $\R\times(0,\infty)$. Then $u$ is an
integral solution and satisfies \eqref{eq:11-entropy} for all smooth
convex entropy pairs. This is conditional on the convergence; it
is not a compactness or existence theorem for general systems.
\end{theorem}
\begin{expandedproof}[moving derivatives onto the compact test]
The chain rule yields
\[
 \eta(u^\varepsilon)_t+Q(u^\varepsilon)_x
 =\varepsilon\eta(u^\varepsilon)_{xx}
 -\varepsilon (u_x^\varepsilon)^TD^2\eta(u^\varepsilon)u_x^\varepsilon.
\]
The last term is nonpositive. Multiply by $\varphi\ge0$ and integrate:
\[
 \iint(\eta(u^\varepsilon)\varphi_t+Q(u^\varepsilon)\varphi_x)
 \ge-\varepsilon\iint\eta(u^\varepsilon)\varphi_{xx}.
\]
The error is $O(\varepsilon)$ because the test support and the range
are compact. Dominated convergence handles the other terms. For the
weak PDE, test the viscous equation directly, including $t=0$:
\[
 \iint(u^\varepsilon\cdot\varphi_t+F(u^\varepsilon)\cdot\varphi_x
 +\varepsilon u^\varepsilon\cdot\varphi_{xx})
 +\int g\cdot\varphi(\cdot,0)=0.
\]
The same argument proves \eqref{eq:11-weak}. No integration over
spatial infinity, energy decay assumption, or gradient bound is
needed for these compactly supported tests. An arbitrary nonsmooth
system entropy is not automatically covered by scalar convolution:
that convolution need not preserve the system flux compatibility.
\end{expandedproof}

\subsection{Scalar entropy and the full contraction argument}
Let $m=1$, $f\in C^2(\R)$, and work in
$u\in C([0,\infty);L^1(\R))\cap L^\infty$ with initial trace $g$.
For each $k\in\R$, the Kruzhkov pair is
\[
 \eta_k(u)=|u-k|,\qquad
 q_k(u)=\operatorname{sgn}(u-k)(f(u)-f(k)).
\]
It follows from the smooth pairs by replacing $|s-k|$ by
$\sqrt{(s-k)^2+\delta^2}$ and taking
$Q_\delta(s)=\int_k^s\eta_\delta'(r)f'(r)\,dr$.
On a bounded range both components converge uniformly, so entropy
inequalities pass to the limit.

\begin{theorem}[Scalar contraction; source \S11.4.3, Theorem 3]
If $u,v$ are scalar entropy solutions in this class with data $g,h$,
then for all $0\le s\le t$,
\begin{equation}\label{eq:11-contraction}
 \|u(\cdot,t)-v(\cdot,t)\|_1
 \le\|u(\cdot,s)-v(\cdot,s)\|_1.
\end{equation}
In particular, the solution with a prescribed initial trace is unique.
No convexity of $f$ is required.
\end{theorem}
\begin{expandedproof}[doubling, diagonal limit, and endpoint times]
For fixed $(y,s)$ use $k=v(y,s)$ in the entropy inequality for $u$,
and for fixed $(x,t)$ use $k=u(x,t)$ in that for $v$. The pair
$q(a,b)=\operatorname{sgn}(a-b)(f(a)-f(b))$ is symmetric in $a,b$.
Add and integrate the inequalities to get, for nonnegative smooth
compactly supported $\zeta(x,t,y,s)$ away from $t=0,s=0$,
\[
 \iiiint \bigl\{|u-v|(\partial_t+\partial_s)\zeta
       +q(u,v)(\partial_x+\partial_y)\zeta\bigr\}\ge0.
\]
Fubini is legitimate because all variables are bounded on the compact
support; the inequalities hold for every constant $k$ and smooth test.
Take
\[
 \zeta=\rho_\delta(x-y)\rho_\delta(t-s)
 \phi\bigl((x+y)/2,(t+s)/2\bigr),\qquad \phi\ge0.
\]
For small $\delta$ the support remains away from both initial times.
The difference-kernel derivatives cancel under the summed derivatives;
the derivatives of $\phi$ have coefficient one.
Use mean and difference variables; their constant positive Jacobian
can be divided out. Translations are continuous in $L^1$ on compact
space-time sets. The functions $|a-b|$ and $q(a,b)$ are jointly
Lipschitz on any bounded square: for fixed $b$, the piecewise derivative
of $q$ is $\pm f'(a)$ and the two pieces meet continuously at $a=b$;
then repeat in $b$. Thus the diagonal limit gives
\begin{equation}\label{eq:11-kato}
 \iint(A\phi_t+B\phi_x)\ge0,\quad
 A=|u-v|,\quad B=q(u,v),\quad |B|\le LA,
\end{equation}
where $L=\max_{|r|\le M}|f'(r)|$ and $M$ bounds both solutions.

Choose $\phi=\alpha_R(x)\beta(\tau)$, with $0\le\alpha_R\le1$,
$\alpha_R=1$ on $[-R,R]$, zero outside $[-2R,2R]$,
$|\alpha_R'|\le C/R$, and $\beta\ge0$ compact in positive time.
The flux contribution tends to zero, since its absolute value is
at most $(CL/R)\|\beta\|_\infty\int\|u-v\|_1d\tau$.
The time integral is finite by $L^1$ continuity on compact intervals.
Dominated convergence yields
$\int\|u-v\|_1\beta'\ge0$. Hence the distributional time derivative
of this continuous function is nonpositive. Equivalently, approximate
the indicator of $[s,t]$ by trapezoidal Lipschitz functions and then
mollify; the two end averages tend to the values at $s,t$ by
continuity. This proves contraction first for $s>0$, then for $s=0$
by the strong initial traces. Continuity also upgrades an a.e.-time
inequality to every time.
\end{expandedproof}
\PDEAnnotation{初始迹与全时刻唯一性}{消黏的分布初值不自动给强
$L^1$ 连续性。本定理把 $C_tL^1_x$ 与强初始迹作为解类条件，
再把几乎处处时刻的收缩提升到每个时刻。一般系统没有相同的序结构。}
The same Kato inequality gives a local finite-speed estimate:
\[
 \int_{|x|\le R}|u(x,t)-v(x,t)|\,dx
 \le\int_{|x|\le R+Lt}|g(x)-h(x)|\,dx.
\]
To verify it, use cutoffs of the backward interval
$|x|<R+L(t-\tau)$ in \eqref{eq:11-kato}; they satisfy
$\phi_\tau+L|\phi_x|\le0$ between the two time cutoffs.
The bound $|B|\le LA$ makes the lateral flux nonpositive.
Smooth approximation and the strong local traces supply the two
end sections. This local version also proves uniqueness for bounded
nonintegrable Riemann data with $L^1_{\rm loc}$ time continuity.
\PDEBodyDiagram{11.4}{The entropy chord for a nonconvex flux}
{For $f(u)=u^3-u$, the chord from $u_l=-1$ to $u_*=1/2$ has slope
$-1/4=f'(1/2)$ and lies below $f$ on that interval. The entropy
solution for right state $1$ joins this shock to a rarefaction
$u=\sqrt{(x/t+1)/3}$ for $-1/4<x/t<2$.
A single jump from $-1$ to $1$ violates the chord criterion.
The exact states and speeds are derived in Problems 14--15.}{ch11-entropy-chord}
